\documentclass[11pt,a4paper]{article}

\usepackage[margin=1in]{geometry}
\usepackage{amsmath, amssymb, amsthm}
\usepackage{mathtools}
\usepackage{bm}
\usepackage{hyperref}
\usepackage{enumitem}
\usepackage{microtype}
\usepackage[T1]{fontenc}
\usepackage{lmodern}
\usepackage[dvipsnames]{xcolor}
\usepackage{booktabs}

\hypersetup{
  colorlinks=true,
  linkcolor=blue!50!black,
  urlcolor=blue!50!black,
  citecolor=blue!50!black
}

\newtheorem{theorem}{Theorem}[section]
\newtheorem*{theorem*}{Theorem}

\title{Structural Foundations of ED-Substrate Gravity:\\
Newton, the Transition Scale, the Combination Rule, and the\\
Baryonic Tully--Fisher Relation}

\author{Allen Proxmire}

\date{April 2026}

\begin{document}

\maketitle

\begin{abstract}
We derive the structural foundations of gravity within the Event Density (ED) substrate framework, working from substrate-native primitives---micro-events, participation channels, relational adjacency, and decoupling surfaces---without importing field-theoretic or thermodynamic machinery. Four principal results are established. First, Newton's law of gravity, $a = GM/R^2$, emerges from a substrate-level cumulative-strain reading of the chain's stability landscape combined with a participation-count bound on decoupling surfaces; the recovery uniquely fixes the substrate UV cutoff to $\ell_{\mathrm{ED}}^2 = \hbar G / c^3$, identifying it with the Planck scale. Second, the gravitational transition acceleration $a_0 = c \cdot H_0 / (2\pi) \approx 1.08 \times 10^{-10}\,\mathrm{m/s^2}$ is derived from a substrate-native dipole mechanism in which the chain's acceleration breaks the three-dimensional isotropy of its participation adjacency, projecting the cosmic decoupling surface onto the chain's accessible region with a leading anisotropic mode whose $2\pi$ azimuthal periodicity sets the effective cosmic rate. The prediction matches the empirically determined MOND acceleration scale within approximately $10\%$, with no free parameters. Third, we ratify substrate rule \textbf{the ED Combination Rule}: in the joint weak-gradient regime, the chain's stability landscape acquires a logarithmic cross-term with coefficient $\sqrt{M\cdot a_0}$, producing the multiplicative effective acceleration $a = \sqrt{a_N \cdot a_0}$. Fourth, composing T19 (Newton), T20 (transition scale), and the ED Combination Rule yields the flat-rotation-curve result $v^2 = \sqrt{G M a_0}$ and the slope-4 baryonic Tully--Fisher relation $v^4 = G M a_0$, with prefactor expressed entirely in substrate constants. The gravitational sector is therefore now \emph{structurally complete and parameter-free}: the full empirical phenomenology of galactic gravity---Newtonian behavior in the high-acceleration regime, the transition at $a_0$, flat rotation curves and slope-4 BTFR in the deep regime---derives from substrate primitives alone, with no tunable parameters.
\end{abstract}

\section{Introduction}\label{sec:intro}

The Event Density (ED) framework treats spacetime, geometry, fields, and forces as emergent rather than fundamental. At its substrate level, ED is built from discrete micro-events, the relational adjacency among them, and stable participation channels through which micro-events propagate. Geometric and field-theoretic structures are coarse-grained summaries of these substrate-level relations.

The present paper derives the structural form of gravitational dynamics---Newton's law in the high-acceleration regime, the transition acceleration that separates the Newtonian from the deep regime, the combination law governing the deep regime, and the flat-rotation-curve / Tully--Fisher phenomenology that follows---from substrate-native primitives alone, without postulating fields, Lagrangians, or thermodynamic ensembles as fundamental.

Three motivations frame this work. \emph{First}, the substrate-level approach offers a framework in which gravity is not assumed but derived from the relational structure of becoming. \emph{Second}, where the derivations succeed, they fix coupling constants and acceleration scales in terms of fundamental ED-substrate quantities, eliminating parameters that field-theoretic approaches treat as free. \emph{Third}, the substrate level reveals structural commitments that the field-theoretic level cannot see---most notably, the multiplicative ED Combination Rule that closes the deep-regime articulation question.

The paper presents four principal structural results.

In Section~\ref{sec:newton}, Newton's gravitational law emerges from a substrate-native cumulative-strain reading of the chain's stability landscape combined with a participation-count constraint on decoupling surfaces. The derivation produces the gravitational constant in substrate quantities: $G = c^3 \ell_{\mathrm{ED}}^2 / \hbar$. Newton-recovery fixes the substrate cutoff at the Planck scale, $\ell_{\mathrm{ED}} = \ell_P$. This result is registered as Theorem T19.

In Sections~\ref{sec:cosmic} and~\ref{sec:dipole}, the transition acceleration is derived. The cosmic horizon, situated at distance $R_H = c/H_0$, contributes to the chain's stability landscape via a substrate-native mechanism whose leading anisotropic mode has $2\pi$ azimuthal periodicity, yielding $a_0 = c \cdot H_0 / (2\pi)$. This result is registered as Theorem T20.

In Section~\ref{sec:ed11}, we ratify the ED Combination Rule: in the joint weak-gradient regime, the chain's stability landscape acquires a logarithmic cross-term with coefficient $\sqrt{M\cdot a_0}$, producing the multiplicative effective acceleration $a = \sqrt{a_N \cdot a_0}$.

In Section~\ref{sec:btfr}, we compose T19, T20, and the ED Combination Rule to derive flat rotation curves $v^2 = \sqrt{G M a_0}$ and the slope-4 baryonic Tully--Fisher relation $v^4 = G M a_0$. The proportionality constant $G \cdot a_0$ is expressed in fundamental substrate quantities, completing the structural derivation of galactic gravitational phenomenology from primitives. This result is registered as Theorem T21.

Section~\ref{sec:empirical} summarizes the empirical content of the now-complete framework, and Section~\ref{sec:conclusion} concludes with the framework's status and forward implications.

\section{ED-Substrate Ontology}\label{sec:ontology}

This section recalls the substrate-level primitives and relations that the subsequent derivations rely on.

\subsection{Micro-events and event density}

At the substrate level, reality consists of discrete micro-events---atomic acts of becoming. The local rate of micro-event production at a region $x$ is the \emph{event density} $\rho(x)$. Higher $\rho$ corresponds to faster local rate of becoming; lower $\rho$ to slower rate. The field $\rho$ is locally smooth at the substrate scale and varies across larger regions.

\subsection{Participation and adjacency}

Micro-events do not exist in isolation. They are linked by \emph{participation}: the relational structure that determines whether one region can integrate the micro-events of another. Participation strength depends on the compatibility of event-density values and gradients between regions. Two micro-events are ``adjacent'' in the substrate-relational sense when they integrate each other's becoming, share participation bandwidth, and maintain coherent relational timing.

\subsection{Stability and chains}

A \emph{chain} is a sequence of micro-events that maintains coherent participation across its extent. A chain's stability at any micro-event step depends on the coherence of its internal micro-event structure across successive steps, the strain produced by mismatch between the chain's local participation environment and its current state, and the integrated environmental gradient encountered by the chain in its accessible region. These three components---coherence, strain, and gradient-strain---together determine the chain's stability landscape, whose gradient experienced by the chain manifests as force.

\subsection{Decoupling surfaces}

A decoupling surface is a participation threshold where reciprocal participation between two regions becomes one-sided. Two decoupling surfaces are central to the present derivations:

\begin{itemize}
\item The \emph{cosmic decoupling surface}, situated at radius $R_H = c/H_0$, beyond which micro-events cannot reach a given observer in finite cosmic time.
\item The \emph{acceleration-induced decoupling surface}, situated at distance $d_a = c^2/a$ behind a chain accelerating at rate $a$, beyond which micro-events emitted from the chain's deceleration direction cannot be integrated by the chain.
\end{itemize}

Both surfaces are 2-spheres in the substrate's emergent geometry, with $4\pi$ solid-angle structure relative to the relevant central observer.

\subsection{The participation-count bound}

A finite causal-domain region admits a finite count of distinguishable participation-channel degrees of freedom on its boundary. Substrate-level UV finiteness forces the existence of a fundamental cutoff length $\ell_{\mathrm{ED}}$. The participation-count bound takes the form
\begin{equation}
N_{\mathrm{dof}} = \frac{A}{\ell_{\mathrm{ED}}^2},
\end{equation}
where $A$ is the boundary area. This is the substrate-native analog of the holographic count.

\subsection{The chain's stability landscape}

A chain at substrate position $x$ with current micro-event state $s$ evaluates its next propagation step against the available local participation environment. The stability score for a candidate next state $e'$ is
\begin{equation}
\Sigma(e') = \mathrm{Coh}(e', s) - \mathrm{Str}(e', \rho_{\mathrm{local}}) - \mathrm{Grad}(e', \nabla\rho),
\end{equation}
where $\mathrm{Coh}$ measures coherence with the current state, $\mathrm{Str}$ measures participation strain, and $\mathrm{Grad}$ measures the strain accumulated against the surrounding ED-gradient field. The chain extends to the available next state with maximum $\Sigma$.

\section{Substrate Derivation of Newton's Law (T19)}\label{sec:newton}

\subsection{The cumulative-strain reading}

The strain term $\mathrm{Str}(e', \rho_{\mathrm{local}})$ in the chain's stability landscape is determined not by the instantaneous local environmental gradient but by the \emph{integrated} environmental gradient from the chain's current location to the chain's natural rest state. For a chain at radius $R$ from mass $M$, the integrated ED-gradient from $R$ to spatial infinity scales as
\begin{equation}
\int_R^\infty \rho_{\mathrm{grad}}(R')\, dR' \propto \frac{M}{R}.
\end{equation}
The chain's effective acceleration is the rate of change of integrated strain with radial position:
\begin{equation}
a \propto \frac{\partial}{\partial R}\!\left(\frac{M}{R}\right) \propto \frac{M}{R^2}.
\end{equation}
This recovers the inverse-square scaling of Newtonian gravity. The proportionality constant remains to be fixed.

\subsection{The participation-count constraint}

Apply the participation-count bound to a holographic 2-sphere at radius $R$ surrounding $M$. The sphere has area $A_R = 4\pi R^2$, and the participation degrees of freedom on it are
\begin{equation}
N_R = \frac{4\pi R^2}{\ell_{\mathrm{ED}}^2}.
\end{equation}
A substrate-level equipartition principle distributes the local mass-energy across these participation degrees of freedom. Setting $\tfrac{1}{2} N_R \cdot k_B T_R = M c^2$ and identifying the local participation rate via an Unruh-analog inversion $T_R = \hbar a / (2\pi k_B c)$, the $2\pi$ factors cancel:
\begin{equation}
a = \frac{M c^3 \ell_{\mathrm{ED}}^2}{R^2 \hbar}.
\end{equation}

\subsection{Identification of the substrate cutoff}

Comparison with Newton's law $a = GM/R^2$ yields
\begin{equation}
G = \frac{c^3 \ell_{\mathrm{ED}}^2}{\hbar},
\end{equation}
equivalently
\begin{equation}
\boxed{\;\ell_{\mathrm{ED}}^2 = \frac{\hbar G}{c^3} = \ell_P^2.\;}
\end{equation}
The substrate UV cutoff is forced by Newton-recovery to be the Planck length. UV-finiteness establishes the existence of $\ell_{\mathrm{ED}}$; Newton-matching fixes its specific value.

\subsection{Result}

The substrate-level rules together yield Newton's law
\begin{equation}
a = \frac{GM}{R^2}, \qquad G = \frac{c^3 \ell_P^2}{\hbar}.
\end{equation}
Newton's gravitational constant is expressed as a relation between fundamental substrate constants. No free parameters are introduced. This result is registered as \textbf{Theorem T19}.

\section{The Cosmic-Horizon Contribution}\label{sec:cosmic}

\subsection{The cosmic decoupling surface}

For any observer, the cosmic decoupling surface is the participation boundary beyond which micro-events cannot reach the observer in finite cosmic time:
\begin{equation}
R_H = \frac{c}{H_0}, \qquad A_H = 4\pi R_H^2 = \frac{4\pi c^2}{H_0^2}, \qquad N_H = \frac{A_H}{\ell_P^2} = \frac{4\pi c^2}{H_0^2 \ell_P^2}.
\end{equation}

\subsection{The cosmic characteristic rate}

The cosmic decoupling surface evolves at rate $H_0$. The associated natural acceleration scale, obtained by converting the cosmic rate via the speed-of-light propagation constraint, is
\begin{equation}
a_{\mathrm{cosmic}} = c \cdot H_0 \approx 6.81 \times 10^{-10}\,\mathrm{m/s^2}.
\end{equation}
This is the bare cosmic-horizon-induced acceleration scale, prior to the geometric projection mechanism developed in Section~\ref{sec:dipole}.

\subsection{The chain's accessible region}

A chain accelerating at rate $a$ has its participation environment shaped by its acceleration-induced decoupling surface at distance $d_a = c^2/a$. In the high-acceleration regime ($d_a \ll R_H$), the chain's local participation environment is dominated by its own decoupling surface. In the low-acceleration regime ($d_a \gg R_H$), the chain's accessible region extends to the cosmic decoupling surface. The crossover occurs at $a_* = c \cdot H_0$.

\section{The $2\pi$ Dipole-Mode Mechanism (T20)}\label{sec:dipole}

\subsection{Acceleration breaks adjacency isotropy}

The participation environment surrounding a chain is three-dimensionally isotropic at rest. When the chain accelerates, its commitment-order direction becomes asymmetric---there is a forward direction (the direction of acceleration) and a backward direction (toward the acceleration-induced decoupling surface). The chain's accessible-region geometry inherits this asymmetry, acquiring a privileged axis along the acceleration direction.

\subsection{Cosmic-horizon projection onto the chain's accessible region}

The cosmic decoupling surface contributes to the chain's stability landscape through the chain's accessible region. Because the chain's accessible region is anisotropic, the cosmic-horizon contribution is not the isotropic full-sphere integration but a directional projection. Expanding in spherical harmonics aligned with the chain's acceleration axis, the leading mode coupling to the chain's anisotropic adjacency is the dipole $(l=1, m=0)$.

\subsection{The $2\pi$ azimuthal periodicity}

The dipole mode $(l=1, m=0)$ is azimuthally symmetric about the chain's acceleration axis. Its only nontrivial spatial periodicity is the $2\pi$ azimuthal cycle. The chain integrates the cosmic-horizon contribution through this mode at effective rate
\begin{equation}
\gamma_{\mathrm{cosmic}}^{\mathrm{eff}} = \frac{H_0}{2\pi}.
\end{equation}
The $2\pi$ reduction is geometric, arising from the azimuthal periodicity of the leading anisotropic projection mode.

\subsection{The transition acceleration}

The chain's experienced fluctuation rate is $\gamma_{\mathrm{chain}} = a/c$. The transition between local-decoupling-surface-dominated dynamics and cosmic-decoupling-surface-dominated dynamics occurs where the rates match:
\begin{equation}
\frac{a}{c} = \frac{H_0}{2\pi} \quad\Longrightarrow\quad \boxed{\;a_0 = \frac{c \cdot H_0}{2\pi}.\;}
\end{equation}
This result is registered as \textbf{Theorem T20}.

\subsection{Numerical comparison}

With $H_0 = 70\,\mathrm{km/s/Mpc} = 2.27 \times 10^{-18}\,\mathrm{s^{-1}}$ and $c = 3.00 \times 10^8\,\mathrm{m/s}$:
\begin{equation}
a_0 \approx 1.08 \times 10^{-10}\,\mathrm{m/s^2}.
\end{equation}
The empirically-determined MOND acceleration constant is $a_0^{\mathrm{emp}} \approx 1.2 \times 10^{-10}\,\mathrm{m/s^2}$. The structural prediction matches the empirical value to within approximately $10\%$, parameter-free.

\subsection{Sensitivity to $H_0$}

\begin{table}[h]
\centering
\begin{tabular}{ccc}
\toprule
$H_0$ (km/s/Mpc) & $a_0$ predicted (m/s$^2$) & Ratio to empirical \\
\midrule
67 & $1.04 \times 10^{-10}$ & 0.86 \\
70 & $1.08 \times 10^{-10}$ & 0.90 \\
73 & $1.13 \times 10^{-10}$ & 0.94 \\
\bottomrule
\end{tabular}
\end{table}

The prediction is robust against the Hubble tension at the level of approximately $15\%$.

\section{The Deep-Regime Combination Rule}\label{sec:ed11}

The derivations of Sections~\ref{sec:newton}--\ref{sec:dipole} establish two structural results: Newton's law in the high-acceleration regime, and the transition acceleration $a_0$. To produce the empirical galactic phenomenology---flat rotation curves and the slope-4 Tully--Fisher relation---a third result is required: the \emph{combination law} governing how local-mass-induced and cosmic-horizon-induced contributions to a chain's stability landscape merge in the deep-acceleration regime. This combination law is supplied by substrate rule \textbf{the ED Combination Rule}.

\subsection{Statement of the Combination Rule}

In the joint weak-gradient regime---where neither local-mass nor cosmic-horizon contribution dominates the chain's accessible region---the chain's stability landscape $\Sigma$ acquires a \emph{logarithmic cross-term} of the form
\begin{equation}
\Sigma_{\mathrm{cross}}(R) = \sqrt{G M a_0} \cdot \log(R / R_0) + \mathrm{const},
\end{equation}
with $R_0$ a substrate-internal reference scale. The coefficient $\sqrt{G\cdot M \cdot a_0}$ reflects multiplicative participation between local-mass and cosmic-horizon contributions to the landscape: it is the geometric mean of the two source-induced strain scales rather than their sum.

The cross-term is \emph{not} a perturbative addition to an additive landscape. It is the substrate's structural response in the regime where both source contributions co-shape the chain's accessible region. In either pure limit ($a_N \gg a_0$ or $a_N \ll a_0$ outside the joint weak-gradient regime), the cross-term either vanishes or reduces to the familiar pure-Newtonian or pure-cosmic forms.

\subsection{Consequence for the effective acceleration}

The chain's experienced acceleration is the gradient of the total stability landscape with respect to radial position. Differentiating the cross-term:
\begin{equation}
\frac{\partial \Sigma_{\mathrm{cross}}}{\partial R} = \frac{\sqrt{G M a_0}}{R}.
\end{equation}
Identifying this with the chain's effective acceleration in the deep-acceleration regime:
\begin{equation}
a = \frac{\sqrt{G M a_0}}{R}.
\end{equation}
Using $a_N = G M / R^2$:
\begin{equation}
\boxed{\;a = \sqrt{a_N \cdot a_0}.\;}
\end{equation}
The effective acceleration in the deep-acceleration regime is the geometric mean of the Newtonian acceleration and the cosmic transition scale. This is the substrate-native multiplicative combination law.

\subsection{Foundational status}

The Combination Rule is registered as a foundational substrate rule. It is neither a perturbative extension of Newton's law nor an empirical fit to galactic data. It is the substrate-level structural commitment governing how multiple source contributions co-shape the chain's stability landscape in the joint weak-gradient regime. The form of the cross-term---logarithmic in $R$, with $\sqrt{M\cdot a_0}$ coefficient---is fixed by the substrate ontology and the requirement of dimensional consistency with the source-induced strain scales established in Sections~\ref{sec:newton} and~\ref{sec:cosmic}.

The Combination Rule supersedes earlier candidate routes to the multiplicative combination law (matter-wave-mediated combinations, Synge-world-function-mediated combinations) by direct substrate articulation.

\section{Flat Rotation Curves and the Slope-4 Baryonic Tully--Fisher Relation (T21)}\label{sec:btfr}

We now compose Theorems T19 and T20 with the ED Combination Rule to derive the empirical galactic phenomenology.

\subsection{Setup}

Consider a chain in circular orbit at radius $R$ from a baryonic mass distribution of total mass $M$, in the deep-acceleration regime where $a_N(R) \ll a_0$. Per the Combination Rule, the chain's effective acceleration is
\begin{equation}
a = \sqrt{a_N \cdot a_0} = \frac{\sqrt{G M a_0}}{R}.
\end{equation}
Centripetal balance requires $a = v^2/R$, where $v$ is the orbital velocity at radius $R$.

\subsection{Flat rotation curves}

Equating centripetal and effective accelerations:
\begin{equation}
\frac{v^2}{R} = \frac{\sqrt{G M a_0}}{R}
\quad\Longrightarrow\quad
\boxed{\;v^2 = \sqrt{G M a_0}.\;}
\end{equation}
The right-hand side is independent of $R$. The orbital velocity asymptotes to a finite, mass-determined value
\begin{equation}
v_{\mathrm{flat}} = (G M a_0)^{1/4}
\end{equation}
throughout the deep-acceleration regime. \emph{This is the structural origin of flat galactic rotation curves at large radii}---derived from substrate primitives, not fit to observations.

\subsection{Slope-4 baryonic Tully--Fisher relation}

Squaring the flat-velocity result:
\begin{equation}
\boxed{\;v_{\mathrm{flat}}^4 = G \cdot M \cdot a_0.\;}
\end{equation}
This is the baryonic Tully--Fisher relation $v_{\mathrm{flat}}^4 \propto M_b$ with slope exactly 4 and prefactor $G\cdot a_0$ expressed in fundamental substrate constants:
\begin{equation}
G \cdot a_0 = \frac{c^3 \ell_P^2}{\hbar} \cdot \frac{c H_0}{2\pi} = \frac{c^4 \ell_P^2 H_0}{2\pi \hbar}.
\end{equation}
No free parameters appear at any step. This result is registered as \textbf{Theorem T21}.

\subsection{The deep-regime $1/R$ force law}

The effective force per unit chain mass in the deep-acceleration regime is $a(R) = \sqrt{G M a_0}/R$, falling off as $1/R$ rather than the Newtonian $1/R^2$. This is the characteristic deep-regime force law that produces flat rotation curves and is the structural signature distinguishing the deep-acceleration regime from the Newtonian one.

\subsection{Radial-acceleration relation}

For arbitrary $(a_N, a_0)$ ratios, composing T19, T20, and the ED Combination Rule yields a single-valued function $a(a_N)$ that interpolates smoothly between the Newtonian limit ($a \to a_N$ for $a_N \gg a_0$) and the deep-regime limit ($a \to \sqrt{a_N \cdot a_0}$ for $a_N \ll a_0$). This is the substrate-level account of the radial-acceleration relation reported empirically by McGaugh, Lelli, and Schombert (2016).

\subsection{Falsifier}

Any galactic system in the certified deep-acceleration regime that significantly departs from $v_{\mathrm{flat}}^4 = G\cdot M \cdot a_0$ falsifies T21 and, depending on which step fails, falsifies one of the upstream commitments T19, T20, or the ED Combination Rule. The relation is sharp, parameter-free, and directly testable against the SPARC catalog and similar surveys.

\section{Empirical Position}\label{sec:empirical}

\subsection{Tested regimes}

\paragraph{Newton regime ($a_N \gg a_0$).} The framework predicts $a = GM/R^2$ with $G = c^3\ell_P^2/\hbar$. This matches the empirical Newtonian gravitational law exactly.

\paragraph{Transition regime ($a_N \approx a_0$).} The framework predicts the transition at $a_0 = c\cdot H_0/(2\pi) \approx 1.08\times 10^{-10}\,\mathrm{m/s^2}$. This matches the empirical MOND constant within approximately $10\%$, parameter-free.

\paragraph{Deep-acceleration regime ($a_N \ll a_0$).} The framework predicts $a = \sqrt{a_N \cdot a_0}$, equivalently $v^2 = \sqrt{G M a_0}$ (flat rotation curves) and $v^4 = G M a_0$ (slope-4 BTFR). The prefactor $G\cdot a_0$ is expressed entirely in substrate constants. This matches the empirical phenomenology of galactic rotation curves and the empirical baryonic Tully--Fisher relation directly.

\subsection{Status}

The framework's empirical position is structurally complete in the gravitational sector. It correctly predicts the existence and value of Newton's gravitational constant; the existence and value of the gravitational transition acceleration; and flat rotation curves with slope-4 BTFR with the prefactor $G\cdot a_0$ derivable in substrate constants. The successes are not phenomenological fits---they are structural derivations from substrate primitives. There are no free parameters in the entire chain from primitives to galactic dynamics.

\subsection{Falsifiers}

\begin{itemize}
\item A galactic system in the deep regime departing from $v^4 = G\cdot M \cdot a_0$ falsifies T21 (and traces back to T19, T20, or the ED Combination Rule).
\item A measured $a_0$ significantly outside the band set by the Hubble tension falsifies T20.
\item A deviation from $G = c^3\ell_P^2/\hbar$ falsifies T19.
\end{itemize}

\section{Conclusion}\label{sec:conclusion}

The substrate-level approach to gravity within the Event Density framework yields four structural results.

\textbf{Newton's law (T19)} emerges from the cumulative-strain reading of the chain's stability landscape combined with the participation-count bound. The gravitational constant is expressed as $G = c^3\ell_P^2/\hbar$, identifying the substrate UV cutoff with the Planck length as a derived rather than postulated relation.

\textbf{The gravitational transition acceleration (T20)} emerges from a substrate-native dipole mechanism: the chain's acceleration breaks the three-dimensional isotropy of its participation adjacency, projecting the cosmic decoupling surface onto the chain's accessible region with a leading dipole mode whose $2\pi$ azimuthal periodicity yields $a_0 = c\cdot H_0/(2\pi)$. The prediction matches the empirical MOND constant to within approximately $10\%$, parameter-free.

\textbf{The deep-regime ED Combination Rule} is ratified as a substrate-native foundational rule. In the joint weak-gradient regime, the chain's stability landscape acquires a logarithmic cross-term with $\sqrt{M\cdot a_0}$ coefficient, producing the multiplicative effective acceleration $a = \sqrt{a_N \cdot a_0}$.

\textbf{The flat-rotation-curve and slope-4 BTFR results (T21)} follow from composing T19, T20, and the ED Combination Rule. Centripetal balance with $a = \sqrt{a_N \cdot a_0}$ yields $v^2 = \sqrt{G\cdot M \cdot a_0}$ (constant in $R$, flat rotation curves) and $v^4 = G\cdot M \cdot a_0$ (slope-4 BTFR with prefactor expressed in substrate constants).

\textbf{The structural position is final.} The gravitational sector of the ED framework is now structurally complete and parameter-free at the substrate level. The full empirical phenomenology of galactic gravity---Newton in the high-acceleration regime, the transition at $a_0$, flat rotation curves and slope-4 BTFR in the deep regime---derives from substrate primitives alone. No free parameters appear in the chain from primitives to galactic dynamics.

The framework's open questions now lie entirely outside the scope of this paper: in the broader gravitational program (Phase-3 Einstein-equation emergence, kernel-parameter inference, the GR-4A speculative quadrant), in cross-arc closure work, and in the empirical-test program at the laboratory and observational level. The substrate-level gravitational scaffolding is complete; downstream work proceeds against a closed foundation.

\begin{thebibliography}{99}

\bibitem{ed06} Proxmire, A. (2026). \emph{Horizons as Event Density Decoupling Surfaces.} ED foundational paper, internal corpus.

\bibitem{ed07} Proxmire, A. (2026). \emph{Event Density and Relativistic Phenomena.} ED foundational paper, internal corpus.

\bibitem{ed10} Proxmire, A. (2026). \emph{Event Density and the Emergence of Spacetime.} ED foundational paper, internal corpus.

\bibitem{edI06} Proxmire, A. (2026). \emph{Fields and Forces in Event Density.} ED foundational paper, internal corpus.

\bibitem{ed11} Proxmire, A. (2026). \emph{the ED Combination Rule: The Deep-Regime Gradient-Combination Rule.} Substrate-native foundational rule, ratified 2026-04-28. \texttt{arcs/arc-SG/ED\_combination\_rule.md}.

\bibitem{arcN} Proxmire, A. (2026). \emph{Theorem N1: V1 Finite-Width Vacuum Kernel.} ED structural-arc result.

\bibitem{arcR} Proxmire, A. (2026). \emph{Arc R: Spin-Statistics, Cl(3,1) Uniqueness, and Dirac Emergence.} ED structural-arc result.

\bibitem{arcQ} Proxmire, A. (2026). \emph{Arc Q: Theorem 17, Gauge-Field-as-Rule-Type.} ED structural-arc result.

\bibitem{arcB} Proxmire, A. (2026). \emph{Arc B: Theorem 18, V1 Kernel Retardation (Kernel-Level Arrow of Time).} ED structural-arc result.

\bibitem{phase3} Proxmire, A. (2026). \emph{Phase 3: Theorem GR1, V1 with Synge World Function.} ED structural-arc result.

\bibitem{mcgaugh2016} McGaugh, S.~S., Lelli, F., \& Schombert, J.~M. (2016). \emph{Radial Acceleration Relation in Rotationally Supported Galaxies.} Phys. Rev. Lett. \textbf{117}, 201101.

\bibitem{mcgaugh2012} McGaugh, S.~S. (2012). \emph{The Baryonic Tully-Fisher Relation of Gas Rich Galaxies as a Test of $\Lambda$CDM and MOND.} Astron. J. \textbf{143}, 40.

\bibitem{milgrom1983} Milgrom, M. (1983). \emph{A modification of the Newtonian dynamics as a possible alternative to the hidden mass hypothesis.} Astrophys. J. \textbf{270}, 365.

\end{thebibliography}

\vfill
\noindent\rule{\textwidth}{0.4pt}\\
\textit{Manuscript prepared April 2026. Gravitational sector closed 2026-04-28 with the ratification of the ED Combination Rule and the derivation of T21.}

\end{document}
